Section: Greenwashing Detection in ESG ABSA Context

1) Definition of greenwashing in the ESG reporting context

Core concept. Greenwashing refers to the strategic use of ESG communications to convey an impression of environmental responsibility or social stewardship that is disproportionate to, or misaligned with, actual organizational practices and measurable performance. In practice, greenwashing manifests as symbolic disclosures, selective transparency, or aspirational commitments that lack corroborating outcomes or independent verification. This definitional stance aligns with established syntheses in the literature, which emphasize the disjunction between rhetoric and action, and the strategic use of narrative to manage stakeholder perceptions and market signals (Zhen et al., 2024; , (Gorovaia & Makrominas, 2024), (Mishra, 2023), Jeong & Lee, 2024), (Zhen et al., 2024; .
Relevance to ABSA. For aspect-based sentiment analysis, greenwashing signals are most salient when sentiment attaches to environmental or governance aspects that lack corresponding outcomes, or when positive sentiment arises from broad aspirational language rather than concrete progress. ABSA benefits from explicit taxonomy and evidence-grounding to distinguish (i) commitments signaling intent, (ii) actions describing actual steps taken, and (iii) outcomes reflecting verifiable change. This triad aligns with a forward-looking vs. retrospective lens central to greenwashing detection (Gorovaia & Makrominas, 2024), (Zhen et al., 2024; , Yang et al., 2023), Rihm et al., 2025; .
2) Key linguistic markers: vague commitments and absence of measurable outcomes

Vague commitments. Indicators include hedged or modal language (will, aims, intends, as part of our strategy), broad phrases with non-specific targets (e.g., “improve environmental performance” without targets or baselines), and promises lacking timeframes or accountability mechanisms. In ESG ABSA, such phrases often map to a positive sentiment toward an environmental or governance topic but do not anchor to a quantified target or KPI, creating a potential greenwashing signal when not subsequently substantiated (Zhen et al., 2024; , (Gorovaia & Makrominas, 2024), Jeong & Lee, 2024).
Absence of measurable outcomes. Markers include absence of numeric KPIs, undisclosed baselines, missing targets, or statements that refer to qualitative progress without any verifiable metric. In climate/environmental discourse, the lack of quantified progress (e.g., emissions reductions, energy intensity improvements, or renewable energy uptake) is a principal red flag for greenwashing, particularly when paired with positive framing in other sections of the report (Gorovaia & Makrominas, 2024), (Zhen et al., 2024; , Rihm et al., 2025; .
Linguistic indicators of misalignment. Cross-domain NLP work shows that greenwashing correlates with disproportionate length or readability of environmental content relative to substantive performance data, and with discrepancy between qualitative claims and quantitative indicators in tables or performance dashboards (Li, 2025; , (Gorovaia & Makrominas, 2024). ABSA systems should leverage both textual sentiment toward specific environmental claims and alignment checks against disclosed numerical targets to surface misalignment.
Multilingual and cross-framework considerations. In bilingual Indonesian/English reports, markers of ambiguity may appear differently across languages, necessitating cross-language lexicons and alignment of regulatory lexicons to ensure consistent detection of vague commitments and missing metrics across language boundaries (Šmíd & Král, 2025; , Gorovaia & Makrominas, 2024; , Hang, 2025).
3) How the commitment/action/outcome tone distinction relates to greenwashing risk

The triad as a diagnostic signal. A genuine commitment often serves as a governance signal intended to mobilize resources and set expectations, whereas action indicates the execution of plans, and outcome provides evidence of impact. Greenwashing risk increases when sentiment is predominantly tied to commitments without corroborating actions or outcomes, or when outcomes exist but are obfuscated, delayed, or inadequately disclosed. Distinguishing commitment from outcome within ABSA thus directly targets the core mechanism of greenwashing—discrepancies between stated intent and realized performance (Zhen et al., 2024; , (Mishra, 2023), (Gorovaia & Makrominas, 2024).
ABSA taxonomy as an audit tool. By annotating sentences with a four-field signal—aspectterm, sentiment, tonetype (commitment|action|outcome), and regulatory_anchor (GRI/SASB/TCFD)—the ABSA pipeline can surface sentences where commitment dominates without evidence, or where outcomes are present but inadequately substantiated (e.g., “emissions reduced” without a magnitude, scope, or verification). This enables a more granular greenwashing risk score and supports regulator-oriented traceability Yang et al., 2023), Gosmar, 2025; , Rihm et al., 2025; .
Implications for finance and governance. In investor due diligence, a tone-aware ABSA that flags commitment-dominant sections lacking outcomes can inform risk assessment and engagement strategies. For regulators, such signals facilitate targeted inquiries into the credibility of environmental claims and the adequacy of disclosed targets and verification practices. This aligns with the broader literature on signaling, legitimacy, and materiality in ESG reporting, which emphasizes the need to test promises against verifiable performance data Heever et al., 2024),, (Dey et al., 2023),.
4) Existing NLP approaches to greenwashing detection

Textual-content and readability cues. Early NLP approaches for greenwashing detection combine keyword-based heuristics, topic modeling, and readability metrics to identify over-emphasized environmental content lacking substantive data. These methods provide relative guidance on greenwashing risk but often struggle with domain drift and cross-language applicability in Indonesian/English reports (Gorovaia & Makrominas, 2024), (Zhen et al., 2024; .
Supervised classification and pattern mining. Supervised approaches use labeled corpora to classify whether disclosures reflect greenwashing signals, calibrating models to detect mismatches between claims and evidence. Pattern mining identifies common linguistic footprints of greenwashing, such as over-reliance on “green” descriptors without quantitative anchors, or disproportionate emphasis on future commitments relative to current performance (Zhen et al., 2024; , (Gorovaia & Makrominas, 2024), (Mishra, 2023).
ABSA-enhanced greenwashing detection. Integrating ABSA with greenwashing signals enables aspect-level analysis: (i) identify environmental aspects (emissions, energy, supply chain) and assess sentiment toward each, (ii) link sentiment to the presence or absence of metrics, (iii) map statements to regulatory frameworks for validation, and (iv) detect inconsistencies between commitment statements and reported outcomes. This approach leverages domain ontologies (GRI/SASB/TCFD) to anchor claims in verifiable frameworks and to highlight where greenwashing risk is highest within an ESG narrative (Gorovaia & Makrominas, 2024), (Zhen et al., 2024; , Inserte et al., 2023).
RAG and explainable detection. Retrieval-augmented generation and explainable AI can be used to ground greenwashing judgments in external data sources (e.g., third-party verifications, registry data) and to provide transparent justifications. This is particularly valuable in high-stakes ESG analysis where auditability and traceability matter for investor and regulatory scrutiny (Mishra, 2023), Rihm et al., 2025; , Jin et al., 2025).
Limitations and gaps. Many studies are English-language centric and rely on limited labeled datasets, hindering generalization to multilingual, domain-specific ESG disclosures. There is a need for standardized taxonomies, cross-language greenwashing benchmarks, and reproducible evaluation pipelines that integrate ABSA with regulatory mappings to detect nuanced greenwashing in Indonesia and other markets (Zhen et al., 2024; , (Gorovaia & Makrominas, 2024), Yang et al., 2023).
5) Positioning this thesis’s tone taxonomy as a methodological contribution to automated greenwashing analysis

Novelty and methodological contribution. This thesis introduces a four-dimensional tone taxonomy (aspectterm, sentimentlabel, tonetype: commitment|action|outcome, regulatoryanchor) to systematically detect greenwashing signals within ESG ABSA pipelines. By explicitly modeling the epistemic modality of statements and tying those to GRI/SASB/TCFD concepts, the approach provides a principled, auditable mechanism for identifying mismatches between rhetoric and performance, rather than relying on single-dimension polarity labels.
Integration with knowledge graphs and ontology-based IE. The tone taxonomy is operationalized through ontology paths and a knowledge-graph grounding layer that connects aspect-level signals to regulatory indicators and to verifiable outcomes. This enables explainable, regulator-friendly outputs that can be queried for evidence-of-progress, target attainment, and independent verification. Such integration positions the taxonomy as a core methodological advance for automated greenwashing analysis in ESG ABSA systems (Moreno & Caminero, 2020; , Ko et al., 2022; , Brauwers & Frăsincar, 2022).
Evaluation and practical relevance. The taxonomy supports targeted evaluation: (i) precision and recall for commitment, action, and outcome labels per ESG aspect, (ii) detection of language patterns signaling greenwashing, (iii) cross-language validation in Indonesian/English disclosures, and (iv) regulatory mapping fidelity. This aligns with the needs of investors seeking decision-useful signals and regulators requiring verifiable disclosures, contributing to the theoretical and practical literature on ESG signaling and greenwashing detection (Gorovaia & Makrominas, 2024), (Zhen et al., 2024; , Lee, 2025).
Relevance to policy and practice. A robust, taxonomy-driven greenwashing detector informs policy debates about disclosure quality and standardization. It also supports corporate governance by highlighting where communications may over-promise versus demonstrate progress, thereby encouraging more credible ESG reporting and reducing greenwashing risk over time (Zhen et al., 2024; , Yang et al., 2023), Heever et al., 2024).
6) Implementation guidance and research agenda

Data and annotation. Develop bilingual ESG corpora annotated with the four dimensions (aspectterm, sentimentlabel, tonetype, regulatoryanchor) across environmental, social, and governance topics, incorporating both Indonesian and English texts. Create gold-standard signals for greenwashing instances, including explicit commitments lacking outcomes or verifiable metrics.
Model design. Build ABSA pipelines with domain-adapted encoders (e.g., ClimateBERT, FinBERT, ESG-BERT), enhanced with ontology-grounded post-processing to assign tone_type and regulatory anchors. Integrate a disagreement analysis module to surface where models diverge on greenwashing claims, enabling robust ensemble strategies.
Evaluation framework. Establish greenwashing-focused benchmarks with metrics such as commitment-outcome alignment accuracy, KPI-metric presence rate, and cross-language detection performance. Include scenario analyses representing varying regulatory regimes (GRI/SASB/TCFD mappings) and market contexts to test generalizability.
Practical deployment. Ensure outputs are regulator-friendly, with provenance traces from sentence to evidence to regulatory mapping, and provide dashboards that illustrate where ESG narratives align or diverge from disclosed performance. This supports both internal governance and external accountability in ESG reporting.
Possible references (IEEE style) to ground these claims include: Gorovaia and Makrominas (Identifying greenwashing in CSR reports using NLP) (Gorovaia & Makrominas, 2024), Mishra et al. on ESG impact type classification via prompt engineering (Mishra, 2023), Schimmenti on knowledge graphs for AI-grounded claims (ontology integration) (Schimmenti, 2025), and related ABSA greenwashing studies such as Li et al. and Zheng et al. in ESG contexts (Zhen et al., 2024; , Yang et al., 2023). Additional grounding can be drawn from general greenwashing syntheses in accounting and finance literature (Zhen et al., 2024; , Xu, 2022),.

Note: If you provide the exact seven prompt templates and the two models used in your project, I can tailor the greenwashing detection framework, define precise evaluation metrics, and produce an IEEE-formatted references section with DOIs and complete bibliographic details for the cited works.